% THESIS ONLY: the overview ADM asked to keep updated. Part 1 (history_ported.tex) is ../sec_exphist.tex,
% i.e. the Overleaf appendix A--M, ported 2026-11-21 (TODO-key cites filled: bash2013limits, hansen2016cma;
% \textit paragraph heads -> subsections; old closing paragraph replaced below because final/ has no LRT).
% Part 2 (this file) adds every experiment after 2026-09-27, each with its README section in a comment.
% Withdrawn metric "frames per extra alarm" is NOT quoted anywhere (README 2.7).
\section{Experiment History} in main.tex (L534-) and the
% per-part drafts 1b, 5, 6, 7, 8, 9, 10 (cut 2026-09-15; recover the long
% versions from git, commit 4f6a1d5, if the thesis needs them).
%
% Budget: ~600 words of prose = ~1.1 IEEE columns, no table, no figure.
% (~750 since the 2026-09-27 imitation paragraph; ~1,250 (~2.2 columns) since the six further
% 2026-09-27 paragraphs, "minimal model" through "timing", ~460 words between them. If long, cut
% those first: timing, then the first noise ablation, then the shaped-noise control.)
% It summarises; experiments that carry a claim are written out in full in
% their own sections, so no number here needs a table to be read.
%
% Carries the corrections README 3.5 lists as owed on the Overleaf parts 2-4:
%   - first experiment had 2 legitimate TX->RX pairs, not 1;
%   - two-agent numbers are job 99211's: BER ~0.24, per-agent power ~0.62,
%     det 2-8%, kurt ~-1.25 (NOT 0.33 / 0.45 / 3-10%); "double the single
%     agent" is gone -- the gain at matched total power is ~15%;
%   - the centralised-optimiser caveat is stated.
% Other numbers from: jobs 101306/101314/101328 (learned detector),
% 101622-101817 (raw-IQ RL), 101860/101866/102115 (characterisation),
% 101870/102305/102316/102319/102390 (realistic channel). The 2026-09-27 paragraphs carry their
% README sections and job IDs in the comment after each.
% ---------------------------------------------------------------------------

\section{Experiment History}
\label{app:history}
We arrived at the current model by extending a simple link one step at a time.
This section summarises what each step established and why we moved on.
Unless stated otherwise, an attacker maximised $R = \mathrm{BER} - \beta D - \gamma P$, with $D$ a
detection penalty and $P$ the jammer power, trained either by PPO or by backpropagating directly
through a differentiable QPSK link built in Sionna.

\subsection{Statistical detectors}
A barrage jammer against two legitimate pairs on a lossless channel validated the chain: BER, power and
detection moved when, and only when, the jammer did.
A PPO jammer with a Gaussian policy then beat a power threshold ($\mathrm{BER} \approx 0.19$ at
$10\%$ detection) by transmitting just below it.
A kurtosis detector stopped it completely, since a Gaussian policy cannot produce the sub-Gaussian
statistics of QPSK, and neither a spline flow nor a Gaussian mixture trained by PPO did better.
The same flow trained by direct gradient evaded at once (kurtosis $-1.25$ against a threshold of
$-1.0$, BER up to $0.17$), but only with a white-box receiver and a genie view of the transmitted
symbol.

\subsection{Two cooperating jammers}
Trained jointly, two jammers reached $\mathrm{BER} \approx 0.24$ at $2$--$8\%$ detection and power
$0.62$ each.
The received constellation converged to four clean QPSK-like clusters with a quarter of the bits
wrong: the team hid by making the sum look legitimate.
At matched total power the gain over one jammer was about $15\%$, and with a single optimiser over
both agents this shows that a coordinated solution exists, not that decentralised agents find it.

\subsection{Learned detection}
The spectrogram CNN of Li \emph{et al.}~\cite{li2022jamming_ofdm} did not replicate on the flat carrier
($78.9\%$ accuracy), and caught barrage noise only at its $2.5\%$ false-alarm rate: there, a jammed
spectrogram is a clean one at lower SNR.
On an 802.11a-like OFDM link it replicated at once ($99.79\%$ accuracy, $0\%$ false alarms).

\subsection{Policy gradient over raw IQ}
MAPPO jammers emitting a complex value per subcarrier never learned against the frozen CNN, whose
output stayed at $0.999$ in every run.
Every broadband action the policy sampled was flagged, so the penalty was the same across the batch
and the normalised advantages carried no signal.
Reducing the action to two dimensions solved stealth but not damage ($\mathrm{BER}$ stuck at $0.35$),
and without the genie symbol cancellation is impossible for i.i.d.\ QPSK.
The obstacle was the action space, not the reward.

\subsection{Detector characterisation}
Sweeping attack parameters through the frozen CNN instead of training against it, interference at
equal power and equal damage ($\mathrm{BER} \approx 0.36$) was flagged $1.1\%$ of the time in-band and
$99.9\%$ once it also touched the empty guard and DC bins.
The replicated detector was a spectral-mask monitor.
Retraining on in-band jamming closed the gap at a price of $3.8\%$ false alarms, and adding a simple
energy detector cut the best stealthy BER from $0.42$ to $0.005$, because on a lossless channel any
added power is visible.

\subsection{Realistic channel}
With per-link frequency-selective fading and zero-forcing equalisation, a sparse in-band jammer imposed
a BER floor of $0.05$--$0.07$ that did not fall with SNR, and fading gave it cover again.
Two of our own numbers did not survive the right comparison.
A $70\%$ gain from channel-aware subcarrier selection, measured at equal configuration, fell to between
$-4.6\%$ and $+6.2\%$ at equal detectability.
And with the detector's own false-alarm rate as the stealth budget, no effective unflagged attacker
existed below $20$\,dB, and the best BER at $20$ and $30$\,dB was $0.011$ and $0.004$.

\subsection{A minimal model with an optimal detector}\label{sec:exphist-minimal}
To simplify, we reduced the link to single QPSK symbols in white Gaussian noise, where the Neyman--Pearson optimal detector has a closed form.
Within its $5\%$ false-alarm rate no attack blind to the symbol caused over $10^{-3}$ excess BER at $0$--$28$\,dB $E_b/N_0$, a finite instance of the square-root law~\cite{bash2013limits} that checked the model rather than adding a finding.
Its value was as a reference: at $8$\,dB, against a blind attack on $2\%$ of the symbols, a CNN on the IQ histogram flagged $12\%$ of frames and the optimal test $85\%$, a gap that shrank as the same power was spread over more symbols.
% README Sec. 3.3 (E1, jobs 2243867/2243879) and 2.1 (prior art, hence "checked the model", not a claim).
% "< 1e-3": largest blind entry in the 3.3 table is gaussian 7e-4 at sigma 0.5 (3.3's prose says
% "3e-4", which is boundary_blind alone). Detector rates re-read from
% artifacts/m0/frontier/results_sigma0.2.json: boundary_blind, P = 0.01, alpha 0.05, duty 0.0196:
% p_l 0.116 / p_np 0.847; duty 1.0: p_l 0.431 / p_np 0.586 (README's "gap 0.73 -> 0.16").
% TODO-key: Bash, Goeckel & Towsley, "Limits of reliable communication with low probability of
% detection on AWGN channels", IEEE JSAC 31(9), 2013 -- bash2013limits in paper_drafts/refs_new.bib,
% not in Overleaf refs.bib (b9620a0).

% CUT for length 2026-09-27 (restore if space; the multi-jammer half is repeated on the waveform link below):
% \textit{Noise and jammer count.}
% On the faded link, noise changed the CNN's false-alarm rate ($0.42$ at $0$\,dB, $0.01$ at $20$\,dB $E_b/N_0$) but not the detection rate of a fixed jammer ($0.42$--$0.60$ throughout).
% More independent jammers at equal total power raised BER and detection together, leaving the median unflagged BER at $0.008$--$0.009$.
% README Sec. 3.3c findings 1 and 4 (sim08_ablation run001; jobs 2261126/2261146/2261173, figures 2261174).
% Fixed jammer: n_active 16, power 1, suite P(det). Finding 4: 1/2/4 jammers, n = 16, mean over
% Eb/N0 >= 15 dB: BER 0.066/0.084/0.096, detection 0.475/0.513/0.544; median confirmed stealthy BER
% 0.0078/0.0092/0.0088, the suite's own clean FAR as the budget.

\subsection{Reproducing a generative jammer}
We reproduced the CGAN jammer of Zhou \emph{et al.}~\cite{11184988} on a single-carrier QPSK waveform link at $30$\,dB SNR.
At a BER of $10^{-3}$ it needed $6.7$\,dB less power than white noise (paper: $2$--$6$\,dB), but with both jammers asynchronous it beat the paper's ``optimal'' same-modulation jammer by $1.3$\,dB rather than trailing it by $1.3$\,dB.
Its losses have no error-rate term, so it did damage only by imitating QPSK.
% README Sec. 3.3b; artifacts/cgan/run001_async_ber_vs_jsr.json (eval job 2260623): noise-GAN 6.66 dB,
% GAN-optimal -1.30 dB vs the paper's +1.31. The withdrawn 2026-09-14 match (4.34 / +1.01, GAN scored
% synchronised) is deliberately not quoted.

\subsection{Classical baselines on the waveform link}
On this link, barrage and pulsed jamming from one to four asynchronous jammers caused no confirmed bit error within the $5\%$ false-alarm rate of one- or two-sided power, kurtosis, or the retrained CNN of~\cite{li2022jamming_ofdm}; only a jammer that knew the symbol did.
The CNN replicated ($99.94\%$ accuracy at $+10$\,dB jamming-to-signal ratio) once spectrograms used a fixed decibel scale: the per-image contrast stretch we had used erases a jammer that lifts the whole band evenly, which likely explains the flat-carrier failure above.
% README Sec. 3.3d (jobs 2266131 CNN, 2266138 K = 1, 2266140 K = 2-4; artifacts/cgan/baselines/run001/
% summary.json). Last sentence: README calls the contrast stretch "a cleaner explanation" of the
% flat-carrier result than the one in "Learned detection" above; the user may prefer to fix it there.
% Cut for length: the exact LRT flags barrage noise ~35 dB below where it does damage (README 3.3d).

\subsection{A learned shaped-noise control}
Gaussian noise with $48$ spectral and temporal gains tuned by CMA-ES~\cite{hansen2016cma} reached a BER of $10^{-3}$ at $15.9$\,dB less power than white noise, by concentrating its power in band and in one of $16$ symbol slots.
Tuned against the CNN, it evaded it only by leaving the victim's band, where it caused no errors.
% README Sec. 3.3e (cgan_learned run001, jobs 2267051/2267090): beta = 0 (BER only) crosses BER 1e-3
% at -16.4 dB vs -0.5 dB for white noise; 99-100 % in-band, effective duty 0.06. CNN-trained: CNN
% 0.03-0.09 up to -20 dB with 0-1 % in-band, BER 0 there. Stealthy BER 0 vs every detector.
% TODO-key: Hansen, "The CMA evolution strategy: a tutorial", arXiv:1604.00772, 2016 -- hansen2016cma
% in paper_drafts/refs_new.bib, not in Overleaf refs.bib (b9620a0).

\subsection{Imitation as initialization}
Our first detector-aware generators followed~\cite{11184988}: the generator was first trained adversarially to imitate the victim's transmit waveform, then trained on the objective of Sec.~\ref{sec:method} for 400 steps.
Both choices were revised.
The budget was short: with ten times the steps, the same generators reached a bit error rate of $3\times10^{-4}$ with up to 3\,dB less power.
The imitation start was a bias toward the victim's modulation.
Trained from random weights against the learned detector, the generator put 97\% of its power on 5\% of the symbols, and at that error rate and 30\,dB SNR the detector flagged $15\%$ of its frames, against $48\%$ for the imitation start (four seeds each).
The price is damage when loud: its error rate levels off near $3\%$, where the imitation-started generator reaches $30\%$, at powers at which every detector already flags every frame.
Imitation alone, with no damage term, needed 12--15\,dB more power for the same error rate.
% README Sec. 3.3f "Imitation start and training budget" + seeds + run002.
% 11184988 = Overleaf key of Zhou 2025 (zhou2025cgan in the other drafts).
% User 2026-09-27: appendix for now; promote to the main text if space.

\subsection{A first noise ablation}
We evaluated the $400$-step generators and classical attacks at $0$--$40$\,dB SNR, recalibrating every detector per level.
For barrage noise, the loudest unflagged power fell by about $0.7$\,dB per dB of SNR against one-sided power and the CNN, while the onset of errors stayed nearly flat above $10$\,dB, so results at the $30$\,dB of~\cite{11184988} do not transfer across SNR.
Its generator gains, measured against the imitation-only generator, are superseded.
% README Sec. 3.3g, the 2026-09-23 run (cgan_snr run001; jobs 2270449/2270556): findings 2-3, stealth
% edge -21 -> -49 dB (power 1-sided), -16 -> -48 dB (CNN) over SNR 0 -> 40; BER onset -2.0 to -3.1 dB
% at SNR >= 10. CNN weights frozen at 30 dB, scale and threshold re-fitted. Its -85.1 pp headline
% compared D2 with D1 and is superseded (README 3.1, 3.3g "E2 on run002"), so it is not quoted.

\subsection{Timing among several jammers}
With the detector at the victim's receiver, a team can at best match one jammer sending its received sum.
For teams of the $400$-step CNN-targeted generator, fading on the jammer links left no loss to recover, but timing mattered: at matched BER and $30$\,dB SNR, four copies with bursts aligned at the receiver matched or slightly beat a single copy at $7.3$\,dB less power each, random offsets cost $6$\,dB and $35$ points of detection probability, and a follower timing error of standard deviation $0.5$--$1$ symbol already lost the match.
% README Sec. 3.3h (E3 pre-check, jobs 2273604-2273606) and 3.3i (D4a, jobs 2273650/2273651). Matched
% excess BER 3e-4, CNN P(det), K = 4, lossless: one 0.822 @ -21.2 dB, aligned 0.649 @ -22.5, random
% 0.996 @ -16.3 (30 dB); D4a knee at 30 dB, and in JSR also at 15 dB. "No loss to recover": fading
% helped a single jammer (block fading acts as frame-level pulsing it could apply itself; 3.3h #1).
% Cut for length: D4a reaches the random-offset floor by sigma ~ 8-16; E3 random timing costs
% 19-35 pp over 15/30 dB (35 is the 30 dB value).
% CANDIDATE: the user may move this paragraph to the main text or to Future Work instead.
% CAUTION: README 3.1 says no E3/team numbers go in the paper yet, and 3.3j says E3/D4a froze the
% under-trained run001/task14_G (400 steps, imitation start) and should be re-run on the current
% generator before being quoted. Included here as history on request; user's call.


\subsection{A learned team policy}
A shared network that gave each of four jammers its power and transmit advance, trained by direct gradient, learned to transmit at full power rather than to compensate the delays between them.
It lost to simply advancing each jammer by its known delay when timing was exact ($0.794$ against $0.705$ detection probability at $30$\,dB) and matched it for timing errors of one symbol or more.
We paused the team direction here.
% README 3.3j (D4b, jobs 2273939/43). Conditional on the under-trained 400-step task14_G.

\subsection{Shadowing on the victim's link}
With per-frame log-normal shadowing of the victim's link at $30$\,dB, an energy detector calibrated without it lost the CNN-targeted generator already at $\sigma=0.1$\,dB and was blind to it at $1$\,dB, while one that knew the gain was unaffected at every level and the retrained CNN stayed the weakest detector.
Shadowing is slow and tracked by the receiver's gain control, so the final experiment models the frame-to-frame variation as small-scale Rician fading instead.
% README 3.3k (S1, jobs 2274945-2274952); 3.4 decision 4.

\subsection{A listening jammer}
Placing the jammer on the receiver's symbol grid and retraining the generators synchronously gave no measurable gain: at a frame error rate of $0.1$ the CNN flagged the CNN-targeted generator ($\beta=1$) on $0.135\pm0.080$ of frames synchronously and $0.128\pm0.022$ asynchronously over three seeds, and a first single-seed gain did not replicate.
% README 3.3l (S2). Synchronous matched QPSK costs 3.6 dB.

\subsection{Training at the operating point}
Retrained at the operating point of Sec.~\ref{sec:method:env}, against a CNN retrained there, the best generators were flagged by the full-band energy detector on about $25\%$ of frames at a frame error rate of $0.1$, against about $63\%$ at $30$\,dB, whether they had been trained at the lower or the higher SNR; the CNN flagged them on $7$--$8\%$.
% README 3.3m (S4, jobs 2275923-2275975). Its generators and defender are the final experiment's regression reference.

\subsection{Noise-level uncertainty}
At the same operating point, a per-frame uncertainty of the noise level of $\sigma_N=2$\,dB raised the false-alarm rate of a power detector calibrated at the nominal level from $0.05$ to $0.23$ ($0.42$ two-sided).
Re-calibrated on frames that carry the uncertainty, so that it keeps its false-alarm rate, it flagged the CNN-targeted generator at a frame error rate of $0.1$ on $0.25$, $0.14$ and $0.09$ of frames for $\sigma_N=0$, $1$ and $2$\,dB.
This is why every detector of the final experiment is calibrated in the environment it is deployed in, and why its energy detector sits after the matched filter, which keeps out the noise outside the signal's band.
% README 3.3n (S5, jobs 2275921-2276029): full-band power, honest CFAR; naive values in brackets there.

\subsection{The defender retrains once}
Retraining the CNN once, on Li et al.'s jammer types with their JSR range widened downward but without the attacker, raised its detection of three independently trained generators at matched damage from $0.16$--$0.18$ to $0.58$--$0.74$ at $30$\,dB, level with the energy detector, at no cost.
At the operating point nothing moved (at most $0.09$), and retraining cost $4$--$5$ points of detection on Li et al.'s own types.
% README 3.3o (D6 round 1). The attacker's next round was parked.

\subsection{The final experiment}
Sections~\ref{sec:method} and~\ref{sec:results} report it: one pipeline in six environments at the operating point, every attacker against every detector, and the generators retrained in each environment.

\subsection{The defender retrains on the grey-box generator}
At the ideal link, retraining the CNN with two seeds of the grey-box CNN-targeted generator as a fifth jammer type cost nothing on Li et al.'s types at $\alpha=0.05$ (pooled detection rate $93.5\%\to93.6\%$) and caught the generator only partly: at an excess frame error rate of $0.5$, from $0.26$ to $0.39$ for a seed it was trained on and from $0.37$ to $0.56$ for a held-out seed.
At the network's own decision rule its false-alarm rate rose from $0.8\%$ to $2.0\%$.
% README 3.3s (jobs 2279746/2279747).

\subsection{Naive baselines at symbol level}
After the final experiment, the study restarted from the simplest attacks on the one-symbol model of Sec.~\ref{sec:exphist-minimal}: QPSK in white noise, with the jammer added per symbol at equal average power.
The simulated symbol error rate matched its closed form at every noise level.
At signal-to-jammer ratios of $0$, $5$, $10$ and $15$\,dB and low noise, the symbol-aware minimum push across the nearest decision boundary~\cite{xie2022adversarial} reached symbol error rates of $1.00$, $0.33$, $0.11$ and $0.033$, against $0.29$, $0.078$, $0.0027$ and $0$ for barrage noise.
A blind push in a random direction did exactly half the damage of the symbol-aware one; a random legitimate symbol did nothing alone above $0$\,dB.
Because the minimum push lands just past the boundary, any receiver noise sends half of the pushed symbols back.
No detector has been applied to this model yet.
% README 3.1 steps 0-1, A.0 (m0/noise_baseline.py, m0/baselines.py, 2026-10-18). These are the results emailed
% to the supervisor after the 2026-10-01 meeting. Xie et al. = xie2022adversarial (added to bib/project.bib 2026-11-21).

\subsection*{What these steps fixed}
These steps fixed the design of Sections~\ref{sec:system}--\ref{sec:method}: a jammer learned by direct gradient through the link and the detector rather than by policy gradient, a hard power budget instead of a power penalty, every detector calibrated to the same false-alarm rate in the environment it is deployed in, comparison at matched damage rather than at matched power or configuration, and an operating point with impairments the defender cannot know.
